\renewcommand{\Anno}{2008}

\renewcommand{\NumeroProblema}{2}

\renewcommand{\Testo}{Marco e Tullio decidono di dividersi una torta a
forma di triangolo equilatero nella maniera seguente:
\begin{itemize}
    \item  Marco sceglie per primo un punto $P$ del triangolo;

	\item Tullio taglia la torta secondo una retta passante per $P$,
	dopo di che sceglie la pi\`u grossa delle due parti in cui la
	torta \`e stata tagliata, lasciando a Marco la parte restante.
\end{itemize}

\begin{enumerate}
    \item Supponiamo che Marco scelga il centro del triangolo.
    Determinare quale percentuale della torta gli toccher\`{a}, come
    minimo.

    \item Supponiamo che Marco scelga un punto diverso dal centro.
    Dimostrare che, se Tullio gioca bene, a Marco rester\`{a} una
    percentuale di torta inferiore al minimo garantito al punto (a).

    \item Estendere l'analisi al caso di torte a forma di triangolo
    qualsiasi.
\end{enumerate}
}

\renewcommand{\Risposta}{\begin{enumerate}
		\item  A Marco toccheranno come minimo i 4/9 della torta.
	
		\item  \ldots
		
		\item Il punto pi\`{u} conveniente per Marco \`{e} il
		baricentro.
	\end{enumerate}}

\renewcommand{\Aiutino}{\begin{enumerate}
	\renewcommand{\labelenumi}{(\alph{enumi})} 
	\item  Se Tullio taglia parallelamente ad un lato, quale
	percentuale della torta tocca a Tullio ed a Marco? Occorre ora mostrare che
	ogni altra retta garantisce a Tullio una percentuale di torta
	inferiore. A tal fine pu\`{o} essere utile considerare la seguente
	figura:
	\begin{center}
		\psset{unit=5ex}
		\pspicture(-2.5,-1)(2.5,2)
	    \pstTriangle[PointSymbol=none](2;90){A}(2;210){B}(2;-30){C}
	    \pstGeonode[PointSymbol=none, PosAngle=100](0,0){P}
	    \pstHomO[PointSymbol=none, PosAngle=150, HomCoef=1 3 div]{B}{A}[B']
	    \pstHomO[PointSymbol=none, PosAngle=150, HomCoef=1 7 div]{B}{A}[B'']
	    \pstHomO[PointSymbol=none, PosAngle=30, HomCoef=1 3 div]{C}{A}[C']
	    \pstTranslation[PointSymbol=none]{A}{B'}{C}[E]
	    \pstInterLL[PointSymbol=none, PosAngle=40]{A}{C}{B''}{P}{C''}
	    \pstInterLL[PointSymbol=none, PosAngle=-45]{B'}{E}{B''}{C''}{D}
	    \psline[fillstyle=solid, fillcolor=blue](P)(D)(B')
	    \psline[fillstyle=solid, fillcolor=blue](P)(C')(C'')
	    \psline(B')(C')
	    \psline(B'')(C'')
	    \endpspicture
	\end{center}

	Attenzione: occorre dimostrare anche che la fetta migliore per
	Tullio \`{e} sempre il quadrilatero.

	\item Dimostrare che, se Marco sceglie un punto diverso dal
	centro, allora Tullio ottiene un risultato migliore del punto
	precedente tagliando secondo la parallela ad un \emph{opportuno}
	lato.
	
	\item Quale punto notevole del triangolo permette di ripetere i
	ragionamenti dei punti precedenti?
\end{enumerate}}

\renewcommand{\SoluzioneUfficiale}{\TitoloSoluzione\DomandaA
Dimostriamo che, se Marco sceglie il centro del triangolo, Tullio
avr\`{a} al massimo i 5/9 della torta (il che avviene se e solo se
taglia secondo una retta parallela ad uno dei lati), e di conseguenza
Marco avr\`{a} come minimo i 4/9 della torta.

Per dimostrarlo consideriamo intanto il caso in cui il taglio di
Tullio \`{e} parallelo ad un lato.  In tal caso, con riferimento alla
Figura 1, avremo che i triangoli $AB'C'$ ed $ABC$ sono simili con
rapporto di similitudine 2/3 (infatti per una ben nota propriet\`{a}
$AP$ \`{e} i 2/3 dell'altezza $AH$), dunque l'area di $AB'C'$ \`{e} i
4/9 dell'area di $ABC$.  Ne segue che Tullio si mangia il trapezio
$BCC'B'$, cio\`{e} 5/9 della torta, lasciando a Marco il triangolo
$AB'C'$, cio\`{e} i 4/9 della torta.
\begin{center}
    \psset{unit=6ex}
    \pspicture(-2.5,-2.3)(2.5,2.5)
    \pstTriangle(2;90){A}(2;210){B}(2;-30){C}
    \pstGeonode[PosAngle={45,-90}](0,0){P}(0,-1){H}
    \pstHomO[PosAngle=150, HomCoef=1 3 div]{B}{A}[B']
    \pstHomO[PosAngle=30, HomCoef=1 3 div]{C}{A}[C']
    \psline(B')(C')
    \psline(A)(H)
    \rput(0,-2){Figura 1}
    \endpspicture
\end{center}

Ci resta solo da dimostrare che ogni taglio non parallelo ad un lato
produce due fette, entrambe inferiori ai 5/9 del totale.  Intanto
questo \`{e} evidente se il taglio passa per un vertice, in quanto
divide il triangolo in due parti uguali.  Se il taglio non passa per
un vertice, allora interseca esattamente due lati.  Supponiamo, senza
perdita di generalit\`{a}, che intersechi $AB$ ed $AC$ in $B''$ e
$C''$, rispettivamente.  Con riferimento alla Figura 2, dobbiamo
pertanto dimostrare che, tra le due fette, quella pi\`{u} grossa \`{e}
sempre quella corrispondente al quadrilatero ($BCC''B''$ nella
figura), e che la sua area \`{e} sempre minore dei 5/9 del totale.
\begin{center}
    \psset{unit=6ex}
    \hfill\pspicture(-2.5,-2.2)(2.5,2.5)
    \pstTriangle(2;90){A}(2;210){B}(2;-30){C}
    \pstGeonode[PosAngle=100](0,0){P}
%     \pstHomO[PosAngle=150, HomCoef=1 3 div]{B}{A}{B'}
    \pstHomO[PosAngle=150, HomCoef=1 7 div]{B}{A}{B''}
%     \pstHomO[PosAngle=30, HomCoef=1 3 div]{C}{A}{C'}
    \pstInterLL[PosAngle=40]{A}{C}{B''}{P}{C''}
%     \psline(B')(C')
    \psline(B'')(C'')
    \rput(0,-1.8){Figura 2}
    \endpspicture\hfill
    \pspicture(-2.5,-2.2)(2.5,2.5)
    \pstTriangle(2;90){A}(2;210){B}(2;-30){C}
    \pstGeonode[PosAngle=100](0,0){P}
    \pstHomO[PosAngle=150, HomCoef=1 3 div]{B}{A}[B']
    \pstHomO[PosAngle=150, HomCoef=1 7 div]{B}{A}[B'']
    \pstHomO[PosAngle=30, HomCoef=1 3 div]{C}{A}[C']
    \pstTranslation[PointSymbol=none]{A}{B'}{C}[E]
    \pstInterLL[PosAngle=40]{A}{C}{B''}{P}{C''}
    \pstInterLL[PosAngle=-45]{B'}{E}{B''}{C''}{D}
    \psline[fillstyle=solid, fillcolor=green](P)(D)(B')
    \psline[fillstyle=solid, fillcolor=green](P)(C')(C'')
    \psline(B')(C')
    \psline(B'')(C'')
    \rput(0,-1.8){Figura 3}
    \endpspicture\hfill\mbox{}
\end{center}


Consideriamo la Figura 3, in cui abbiamo anche tracciato la parallela
ad $AC$ passante per $B'$.  I triangoli $PC'C''$ e $PB'D$ sono
congruenti in quanto $PC'=PB'$, gli angoli in $P$ sono uguali in
quanto opposti al vertice, e $\angle PC'C''=\angle PB'D=60^{\circ}$.
Ne segue che
$$\mbox{Area}(PC'C'')=\mbox{Area}(PB'D)<\mbox{Area}(PB'B''),$$
e pertanto
\begin{eqnarray*}
    \mbox{Area}(BCC''B'') & = &
\mbox{Area}(BCC'PB'')+\mbox{Area}(PC'C'')  \\
     & < & \mbox{Area}(BCC'PB'')+\mbox{Area}(PB'B'')  \\
     & = & \mbox{Area}(BCC'B')=\frac{5}{9}\mbox{Area}(ABC).
\end{eqnarray*}

In realt\`{a} abbiamo cos\`{\i} dimostrato che 
\begin{eqnarray*}
	\mbox{Area}(BCC'B')-\mbox{Area}(BCC''B'') & = &
	\mbox{Area}(PB'B'')-\mbox{Area}(PC'C'')  \\
	 & = & \mbox{Area}(B'B''D).
\end{eqnarray*}

Poich\'{e} la retta $B'D$ non dipende dalla posizione del punto $B''$,
\`{e} facile osservare che l'area di $B'B''D$ \`{e} tanto maggiore
quanto pi\`{u} il punto $B''$ \`{e} vicino a $B$.

Possiamo pertanto concludere che l'area del quadrilatero $BCC''B''$
decresce man mano che $B''$ si sposta da $B'$ a $B$, variando da un
massimo di 5/9 (corrispondente a $B''=B'$) ad un minimo di 1/2
(corrispondente a $B''=B$) dell'area di $ABC$.  Quindi in ogni caso
tale quadrilatero rappresenta la fetta pi\`{u} grossa.

\DomandaB Tracciamo le tre rette parallele ai lati e passanti per il
centro, ed indichiamo le intersezioni con i lati come nella Figura 4.

Supponiamo che Marco scelga un punto $P$ diverso dal centro del
triangolo.  In tal caso \`{e} facile verificare che $P$ sta
all'interno (e non sul bordo) di \emph{almeno} uno dei triangoli
$AB'C''$, $BC'A''$, $CA'B''$.  Supponiamo, senza perdita di
generalit\`{a}, che $P$ stia in $AB'C''$.  Allora Tullio, tagliando
secondo la parallela a $BC$ passante per $P$, ottiene un trapezio di
area maggiore dei 5/9 dell'area di $ABC$, lasciando dunque a Marco
meno dei 4/9 che avrebbe ottenuto di sicuro scegliendo come $P$ il
centro del triangolo.
\begin{center}
    \hfill
    \psset{unit=6ex}
    \pspicture(-2.5,-2.5)(2.5,2.5)
    \pstTriangle(2;90){A}(2;210){B}(2;-30){C}
    \pstHomO[PosAngle=150, HomCoef=1 3 div]{B}{A}[C'']
    \pstHomO[PosAngle=150, HomCoef=2 3 div]{B}{A}[C']
    \pstHomO[PosAngle=270, HomCoef=1 3 div]{C}{B}[A'']
    \pstHomO[PosAngle=270, HomCoef=2 3 div]{C}{B}[A']
    \pstHomO[PosAngle=30, HomCoef=1 3 div]{A}{C}[B'']
    \pstHomO[PosAngle=30, HomCoef=2 3 div]{A}{C}[B']
    \psline(A')(B'')
    \psline(B')(C'')
    \psline(C')(A'')
    \rput(0,-2){Figura 4}
    \endpspicture
    \hfill
    \pspicture(-2.5,-2.5)(2.5,2.5)
    \pstTriangle(2;90){A}(2;210){B}(2;-30){C}
    \pstGeonode[PosAngle=45](0.4;37){P}
    \pstHomO[PosAngle=150, HomCoef=1 3 div]{B}{A}[C'']
    \pstHomO[PosAngle=150, HomCoef=2 3 div]{B}{A}[C']
    \pstHomO[PosAngle=270, HomCoef=1 3 div]{C}{B}[A'']
    \pstHomO[PosAngle=270, HomCoef=2 3 div]{C}{B}[A']
    \pstHomO[PosAngle=30, HomCoef=1 3 div]{A}{C}[B'']
    \pstHomO[PosAngle=30, HomCoef=2 3 div]{A}{C}[B']
    \pstTranslation[PointSymbol=none]{B'}{P}{C''}[E]
    \pstInterLL[PointSymbol=none]{A}{B}{E}{P}{F}
    \pstInterLL[PointSymbol=none]{A}{C}{E}{P}{F'}
    \psline[fillstyle=solid, fillcolor=green](B)(C)(F')(F)(B)
    \psline(F)(F')
    \psline[linestyle=dashed](A')(B'')
    \psline[linestyle=dashed](B')(C'')
    \psline[linestyle=dashed](C')(A'')
    \psdot(P)
    \rput(0,-2){Figura 5}
    \endpspicture
    \hfill\mbox{}
\end{center}

\DomandaC Nel caso di un triangolo generico il punto pi\`{u}
conveniente per Marco \`{e} il baricentro.  Infatti tutti i
ragionamenti precedenti si possono ripetere parola per parola pur di
sostituire il centro con il baricentro.  In particolare sono ancora
vere le seguenti affermazioni.
\begin{itemize}
    \item Se Marco fissa il punto $P$ nel baricentro, e Tullio taglia
    secondo una retta parallela ad un lato, allora Tullio ottiene i
    5/9 della torta e lascia i 4/9 a Marco (vedi Figura $1'$).  Il
    punto essenziale \`{e} ancora che i triangoli $AB'C'$ ed $ABC$
    sono simili con rapporto di similitudine 2/3 (per una ben nota
    propriet\`{a} del baricentro).

    \item Se Marco fissa il punto $P$ nel baricentro, e Tullio taglia
    secondo una retta non parallela ai lati, allora la parte pi\`{u}
    grossa \`{e} sempre il quadrilatero, ma la sua area \`{e} comunque
    meno dei 5/9 del totale (vedi Figura $3'$).  La dimostrazione si
    basa ancora sulla congruenza dei triangoli $PC'C''$ e $PB'D$, dove
    $D$ \`{e} stato ottenuto tracciando per $B'$ la parallela ad $AC$
    (l'uguaglianza $PC'=PB'$ segue dal fatto che $P$ sta sulla mediana
    relativa al lato $BC$).

\begin{center}
    \hfill
    \psset{unit=6ex}
    \pspicture(-2.5,-2.3)(2.5,2.5)
    \pstTriangle(0,2){A}(-1.5,-1){B}(2.5,-1){C}
    \pstCGravABC[PosAngle=45]{A}{B}{C}{P}
    \pstInterLL[PosAngle=-90]{A}{P}{B}{C}{M}
    \pstHomO[PosAngle=150, HomCoef=1 3 div]{B}{A}[B']
    \pstHomO[PosAngle=30, HomCoef=1 3 div]{C}{A}[C']
    \psline(B')(C')
    \psline(A)(M)
    \rput(0,-2){Figura $1'$}
    \endpspicture\hfill
    \pspicture(-2.5,-2.2)(2.5,2.5)
    \pstTriangle(0,2){A}(-1.5,-1){B}(2.5,-1){C}
    \pstCGravABC[PosAngle=100]{A}{B}{C}{P}
    \pstHomO[PosAngle=150, HomCoef=1 3 div]{B}{A}[B']
    \pstHomO[PosAngle=150, HomCoef=1 7 div]{B}{A}[B'']
    \pstHomO[PosAngle=30, HomCoef=1 3 div]{C}{A}[C']
    \pstTranslation[PointSymbol=none]{A}{B'}{C}[E]
    \pstInterLL[PosAngle=40]{A}{C}{B''}{P}{C''}
    \pstInterLL[PosAngle=-45]{B'}{E}{B''}{C''}{D}
    \psline[fillstyle=solid, fillcolor=green](P)(D)(B')
    \psline[fillstyle=solid, fillcolor=green](P)(C')(C'')
    \psline(B')(C')
    \psline(B'')(C'')
    \rput(0,-1.8){Figura $3'$}
    \endpspicture\hfill\mbox{}
\end{center}

    \item Se Marco sceglie un punto $P$ diverso dal baricentro, allora
    $P$ sta nell'interno (e non sul bordo) di almeno uno dei triangoli
    $AB'C''$, $BC'A''$, $CA'B''$ (vedi Figura $4'$).  Se, per esempio,
    $P$ sta nell'interno di $AB'C''$ (Figura $5'$) Tullio riesce,
    tracciando la parallela a $BC$, a garantirsi pi\`{u} dei 5/9 della
    torta.
\end{itemize}
\begin{center}
    \hfill
    \psset{unit=6ex}
    \pspicture(-2.5,-2.5)(2.5,2.5)
    \pstTriangle(0,2){A}(-1.5,-1){B}(2.5,-1){C}
    \pstCGravABC[PosAngle=90]{A}{B}{C}{G}
    \pstHomO[PosAngle=150, HomCoef=1 3 div]{B}{A}[C'']
    \pstHomO[PosAngle=150, HomCoef=2 3 div]{B}{A}[C']
    \pstHomO[PosAngle=270, HomCoef=1 3 div]{C}{B}[A'']
    \pstHomO[PosAngle=270, HomCoef=2 3 div]{C}{B}[A']
    \pstHomO[PosAngle=30, HomCoef=1 3 div]{A}{C}[B'']
    \pstHomO[PosAngle=30, HomCoef=2 3 div]{A}{C}[B']
    \psline(A')(B'')
    \psline(B')(C'')
    \psline(C')(A'')
    \rput(0,-2){Figura $4'$}
    \endpspicture
    \hfill
    \pspicture(-2.5,-2.5)(2.5,2.5)
    \pstTriangle(0,2){A}(-1.5,-1){B}(2.5,-1){C}
    \pstGeonode[PosAngle=90](0.4;37){P}
	\pstHomO[PosAngle=150, HomCoef=1 3 div]{B}{A}[C'']
	\pstHomO[PosAngle=150, HomCoef=2 3 div]{B}{A}[C']
	\pstHomO[PosAngle=270, HomCoef=1 3 div]{C}{B}[A'']
	\pstHomO[PosAngle=270, HomCoef=2 3 div]{C}{B}[A']
	\pstHomO[PosAngle=30, HomCoef=1 3 div]{A}{C}[B'']
	\pstHomO[PosAngle=30, HomCoef=2 3 div]{A}{C}[B']
    \pstTranslation[PointSymbol=none]{B'}{P}{C''}[E]
    \pstInterLL[PointSymbol=none]{A}{B}{E}{P}{F}
    \pstInterLL[PointSymbol=none]{A}{C}{E}{P}{F'}
    \psline[fillstyle=solid, fillcolor=green](B)(C)(F')(F)(B)
    \psline(F)(F')
    \psline[linestyle=dashed](A')(B'')
    \psline[linestyle=dashed](B')(C'')
    \psline[linestyle=dashed](C')(A'')
    \psdot(P)
    \rput(0,-2){Figura $5'$}
    \endpspicture
    \hfill\mbox{}
\end{center}
}

\renewcommand{\SoluzioniAlternative}{\TitoloSoluzioneAlternativaA Il
punto fondamentale per la soluzione della prima parte del problema, e
cio\`{e} il fatto che l'area del trapezio $BCC''B''$ cresce man mano
che $B''$ si sposta da $B'$ verso $B$, si pu\`{o} dimostrare anche per
via computazionale.  Con riferimento alla Figura~3, indichiamo con
$\theta$ l'angolo $\angle C'PC''$, che pu\`{o} variare da $0^{\circ}$
a $30^{\circ}$. Si ha allora che
$$\mbox{Area}(PC'C'')=\frac{1}{2}\left[PC'\right]^{2}
\frac{\sin\theta\cdot\sin 60^{\circ}}{\sin(60^{\circ}+\theta)},$$
$$\mbox{Area}(PC'C'')=\frac{1}{2}\left[PB'\right]^{2}
\frac{\sin\theta\cdot\sin 120^{\circ}}{\sin(120^{\circ}+\theta)}=
\frac{1}{2}\left[PC'\right]^{2}
\frac{\sin\theta\cdot\sin 60^{\circ}}{\sin(60^{\circ}-\theta)},$$
e pertanto
\begin{eqnarray*}
	\mbox{Area}(BCC'B')-\mbox{Area}(BCC''B'') & = &
	\mbox{Area}(PB'B'')-\mbox{Area}(PC'C'')  \\
	 & = & \frac{\sqrt{3}}{4}\left[PC'\right]^{2}\sin\theta\cdot
	 \frac{\sin(60^{\circ}+\theta)-\sin(60^{\circ}-\theta)}{
	 \sin(60^{\circ}-\theta)\cdot\sin(60^{\circ}+\theta)}.
\end{eqnarray*}

Tale espressione pu\`{o} essere studiata ricorrendo al calcolo
differenziale oppure, pi\`{u} semplicemente, utilizzando le formule
trigonometriche
$$\sin(60^{\circ}+\theta)-\sin(60^{\circ}-\theta)=
2\cos 60^{\circ}\cdot\sin\theta=\sin\theta,$$
$$\sin(60^{\circ}-\theta)\cdot\sin(60^{\circ}+\theta)=
\cos(2\theta)-\cos 120^{\circ}=\frac{3-4\sin^{2}\theta}{2}.$$

Ne segue che
$$\mbox{Area}(BCC'B')-\mbox{Area}(BCC''B'')=
\frac{\sqrt{3}}{2}\left[PC'\right]^{2}\frac{\sin^{2}\theta}{
3-4\sin^{2}\theta}.$$

Quando $\theta$ varia da $0^{\circ}$ a $30^{\circ}$ il numeratore
aumenta ed il denominatore diminuisce (rimanendo sempre positivo),
dunque il rapporto cresce.
}

\renewcommand{\Osservazioni}{
}

\renewcommand{\Commenti}{
}